% Chapter 2: Literature Review
% Dense survey (~1.5 pages) motivating Ch.3-9

\chapter{Literature Review}
\label{ch:literature-review}

This dissertation builds on three intersecting research streams: single-agent policy gradient methods, multi-agent reinforcement learning with opponent modeling, and meta-learning for continuous adaptation. We review the foundational work in each area, with particular attention to the algorithmic lineage that culminates in Meta-MAPG~\cite{kim2021policygradientalgorithmlearning}, the core theoretical object of this dissertation. The final section connects these ideas to the emerging challenge of steering autonomous language model agents.

\section{Single-Agent Policy Gradient Methods}
\label{sec:lit-single-agent}

Policy gradient methods optimise a parameterised policy by estimating the gradient of expected return with respect to policy parameters. \citet{williams1992} introduced the REINFORCE algorithm, which uses the score function (log-derivative) estimator to compute unbiased gradients from sampled trajectories, establishing the Monte Carlo policy gradient framework. \citet{sutton2018rl} provide the canonical exposition, deriving the Policy Gradient Theorem (Theorem~13.1 therein, proved in Chapter~\ref{ch:pgt-proof} of this dissertation), which shows that the gradient may be computed without differentiating the on-policy state distribution---a critical simplification that underpins all modern policy gradient methods. \citet{schulman2015trpo} and \citet{schulman2017ppo} introduced trust-region constraints (TRPO) and clipped surrogate objectives (PPO), respectively, addressing the instability of large policy updates by restricting the policy change per iteration; these methods dominate contemporary single-agent RL practice and motivate the gradient damping considerations in Chapter~\ref{ch:convergence}.

\section{Multi-Agent Reinforcement Learning}
\label{sec:lit-marl}

\citet{littman1994} formalised multi-agent RL via \emph{Markov games} (also called stochastic games, due to \citet{shapley1953}), extending the Markov decision process to environments where multiple agents' policies jointly determine the transition dynamics. This introduces the core difficulty of multi-agent learning: \emph{non-stationarity}---each agent's environment is a moving target, shifting as other agents learn. \citet{busoniu2008} survey early approaches, most of which treat other agents as part of a stationary environment, ignoring the coupling between learners. \citet{lowe2017} propose multi-agent deep deterministic policy gradient (MADDPG), which uses centralised training with decentralised execution: each agent observes others' actions during training to stabilise learning, but acts independently at test time. MADDPG and related methods improve sample efficiency but do not explicitly model how an agent's policy influences the other agents' learning dynamics---the gap that motivates this dissertation.

\section{Learning Awareness and Opponent Modeling}
\label{sec:lit-opponent-modeling}

\citet{zhang2010} introduce \emph{policy prediction}, where an agent maintains an explicit model of the opponent's policy and predicts its next action; this improves performance in competitive games but assumes a fixed opponent model. \citet{foerster2018lola} take a critical step forward with Learning with Opponent-Learning Awareness (LOLA): agent~$i$ optimises against agent~$j$'s policy \emph{after} $j$ takes a gradient step, using a first-order Taylor approximation to anticipate $j$'s update. This captures the \emph{peer-learning} channel---how agent~$i$'s policy shapes agent~$j$'s next policy via the shared training data. LOLA is limited to one-step lookahead and two agents; \citet{letcher2019} extend it to higher-order derivatives (Stable Opponent Shaping, SOS) and analyse convergence in differentiable games. However, LOLA and SOS ignore the agent's \emph{own} multi-step learning dynamics: they optimise the first step but do not account for how the agent's initial policy shapes its own future adaptation. Chapter~\ref{ch:meta-learning} formalises this gap and shows how Meta-MAPG fills it.

\section{Meta-Learning for Multi-Agent Systems}
\label{sec:lit-meta-learning}

Meta-learning, or ``learning to learn,'' optimises for rapid adaptation rather than performance on a single task. \citet{finn2017} introduce Model-Agnostic Meta-Learning (MAML), which finds initial parameters that yield fast gradient-based adaptation across a distribution of tasks; MAML has become the dominant paradigm for few-shot RL. \citet{stadie2018} apply meta-learning to exploration, using the meta-objective to bias the agent toward policies that facilitate future learning. \citet{alshedivat2018} extend MAML to competitive and non-stationary environments, proposing \emph{continuous adaptation via meta-learning}: the agent optimises its initial policy to maximise performance after a sequence of gradient updates. This accounts for the agent's own learning trajectory (the ``own learning'' channel) but assumes other agents' parameters are independent---the complementary limitation to LOLA. \citet{kim2021policygradientalgorithmlearning} unify these threads with the Meta-Multiagent Policy Gradient theorem (Theorem~\ref{thm:meta-mapg}), which computes the exact gradient with respect to both the own-learning and peer-learning channels. The proof (Chapter~\ref{ch:meta-mapg}) and convergence analysis (Chapter~\ref{ch:convergence}) form the theoretical core of this dissertation.

\section{Applications to Language Model Steering}
\label{sec:lit-llm-steering}

Recent large language models exhibit emergent multi-agent phenomena when deployed in interactive or autonomous settings. \citet{christiano2017} and \citet{ouyang2022} use reinforcement learning from human feedback (RLHF) to align LLM outputs with human preferences, treating the language model as a policy and human ratings as a reward signal; this is single-agent RL in token space. \citet{wei2022cot} show that chain-of-thought prompting improves reasoning by conditioning the model on intermediate steps, introducing a sequential decision structure that resembles a policy over reasoning trajectories. \citet{bai2022constitutional} propose constitutional AI, which uses AI-generated feedback to iteratively refine model behaviour---analogous to a meta-learning loop where the model's outputs inform its own training signal. When multiple LLM agents interact autonomously, as in~\cite{agentsofchaos2026}, the environment becomes a stochastic game: each agent's prompts and actions alter the shared state (e.g., a messaging server), creating feedback loops invisible to single-agent methods. Chapter~\ref{ch:llm-steering} formulates LLM steering as a multi-agent policy gradient problem, and Chapters~\ref{ch:cooperative-game} and~\ref{ch:simulations} test whether Meta-MAPG can stabilise joint learning in simulated multi-agent text environments. The central hypothesis is that accounting for peer-learning dynamics (Term~3 of Theorem~\ref{thm:meta-mapg}) prevents the emergent failure modes documented in~\cite{agentsofchaos2026}.
